\section{Conceptual Framework}
\xinze{XXX Unfinished I am trying to simply explain our terminologies here. I have a more explicit version.....XXX. We may delete these or put these somewhere else later}

The evolution of mathematical practice can be understood through two co-constitutive dimensions, \textbf{inherited mathematical structure} and \textbf{human organizational process}. The former is the product of the latter after sedimentation under logical constraints, and the latter is the reorganization of the former within new tooling environments. Each generation's organizational process eventually sediments into the next generation's inherited structure, and each generation's inherited structure in turn determines the starting point and boundaries of its organizational process.

\begin{finding}
    \textbf{Inherited mathematical structure} is the residue of past \underline{human organizational processes} after sedimentation (being recorded and preserved in media), selection (errors discovered and eliminated, ineffective formulations forgotten, inconsistencies corrected), and solidification (becoming "standard," "natural," "obvious").

    It is the starting point that the mathematical community, through historical accumulation, hands to each generation of working mathematicians: the hierarchy of definitions, the dependencies among theorems, the paths of proofs, the division into branches, what counts as "basic," what counts as "advanced." Contemporary mathematicians do not start from zero — they inherit an entire conceptual system that has already undergone logical testing and historical selection. This system exists in the form of textbooks, papers, curricula, and formalized libraries, and is internalized as mathematical intuition through education and practice.
\end{finding}
\begin{finding}
    \textbf{Human organizational process} is the reorganization of \underline{inherited mathematical structure} — partitioning, packaging, presenting, simplifying, generalizing — by mathematicians (and increasingly, machine participants) within the tooling environment of a given historical period, as well as the ongoing development and new creation built upon existing theory under those historical conditions: new definitions, new theorems, new proof paths, new branches.
    
    It is the "living" side of mathematical practice: mathematicians of a given era, carrying intuitions shaped by history, repartition, repackage, and re-present the mathematical content they have inherited, working within the constraints of their contemporary tools — while simultaneously advancing new definitions, new theorems, and new proof strategies on that inherited foundation. This process in turn shapes the intuitions of the next generation: what organizational patterns "feel right," what structures are "useful," what directions are "worth pursuing."
\end{finding}

The central insight of this analytical framework is that \module{Mathlib} is a \textbf{transitional artifact}. It is the first large-scale effort to migrate traditional mathematical knowledge — developed within the old tooling environment of paper, journals, and limited physical space — into a new environment of formalization, Git, and internet-based collaboration. As such, it carries the traces of both eras simultaneously:

Its logical structure (premise dependencies) encodes centuries of accumulated inherited mathematical structure — but these structures have already been reshaped by contemporary tools (Lean's type theory, the tactic system). Between \textbf{what} is carried and \textbf{how} it is carried, inherited structure dominates the content (the theorems are still those theorems; the broad framework of dependencies is unchanged), while contemporary tooling dominates the form (the specific dependency chains, proof paths, and type-theoretic representations are reshaped in the process of migration). This makes the premise dependency network neither a direct copy of traditional mathematics nor an entirely new creation from scratch, but a \textbf{tool-mediated migration product}.

Its organizational structure (module partitioning, namespaces) reflects the constraints and choices of the contemporary tooling environment — yet many of these organizational decisions are continuations of older habits. Top-level modules such as \module{Mathlib.Analysis}, \module{Mathlib.Topology}, and \module{Mathlib.Algebra} correspond directly to traditional disciplinary divisions, and the inertia of these divisions has persisted across the change in tooling environment. At the same time, contemporary tools have given rise to organizational categories that have no counterpart in traditional mathematics (e.g., \module{Mathlib.Tactic}, \module{Mathlib.Lean}) — these are direct products of the new tooling environment. Between the continuation of old habits and the shaping force of new tools, old habits still dominate most of the organizational structure, but new tools are creating entirely novel organizational categories at the margins — margins that may foreshadow future modes of organization.

The \textbf{tension} between the two — the incomplete alignment between inherited structure and contemporary organization — is precisely the most valuable feature of this transitional period for research. When traditional mathematical concept-relations are migrated into the formalized environment, they may conflict with the module boundaries driven by contemporary tooling: logically tightly coupled declarations may be split across different modules, or logically independent content may be bundled under the same namespace due to historical inertia. These misalignments are the observable traces of the collision between the two forces.

This makes \module{Mathlib} simultaneously a mathematical object (whose premise dependency network encodes the logical structure of mathematical knowledge) and a socio-technical object (whose modular organization encodes the practices of the contemporary mathematical community within a specific tooling environment). The separation, entanglement, and conflict between the two constitute the core problem space of our research.